\documentclass[11pt,letterpaper]{article}

\usepackage[margin=1in]{geometry}
\usepackage[T1]{fontenc}
\usepackage[utf8]{inputenc}
\usepackage{times}
\usepackage{microtype}

% Business-letter layout: blocks separated by space, no first-line indent, ragged
% bottom so the signature block is not stretched down the page.
\setlength{\parindent}{0pt}
\setlength{\parskip}{\medskipamount}
\raggedbottom
\pagestyle{empty}

\begin{document}

6 October 2026

\bigskip

Dear Editors,

We submit the manuscript ``GVS5H: Zero-Shot Self-Orchestration with Ledger-Based Control
Improves Coding in Language Models'' for consideration in \textit{IEEE Transactions on Cybernetics}.

Scaffolded and multi-agent language model systems are now standard practice across research and
industry, on the premise that decomposing a problem improves on a single long response. This
premise is seldom evaluated against a matched control, namely the same model given the same
prompt in one call. The manuscript reports that comparison for eleven open- and closed-weight
models on the 100 most recent hard problems in LiveCodeBench. The method requires no training
and no task-specific prompting, and consists of one model called repeatedly in a fresh context,
coordinating with itself through a small set of files on disk.

Every model evaluated in both arms improved, by as much as 25.6 points. Several smaller
open-weight models, once orchestrated, match or exceed the accuracy of much larger proprietary
models at a small fraction of their cost. These results indicate that inference-time
organization can substitute for model scale, which bears directly on the cost and the
accessibility of capable code generation. We acknowledge that the improvement is not universal,
and several configurations show no change or a decline.

The manuscript addresses a question the field has not settled, which is how much of the reported
benefit of multi-agent systems comes from coordination itself rather than from the models or the
training behind it. The answer matters beyond the immediate subfield. Frontier-level coding ability
has so far been available mainly through large proprietary models, and these results
show that comparable accuracy can be obtained from smaller open-weight models that a single
laboratory can host, which affects who is able to use and to study such systems. We believe this
suits the broad readership of \textit{IEEE Transactions on Cybernetics}.

The code needed to reproduce our results is available in our public GitHub repository. Since its
release, we have been encouraged by the interest it has received, including on social media,
which we hope reflects a broader interest in methods of this kind.

On behalf of all authors, thank you for your consideration.

\bigskip

Sincerely,

\bigskip

Simon Lee

Persis Capital Inc., Toronto, Ontario, Canada

slee@persisholdings.com

\end{document}
